% Kapitel "Data Pipeline, Baseline und Experiment-Tracking" (Jan-David Wiederstein).
% Eigenständig einbindbar: % Kapitel "Data Pipeline, Baseline und Experiment-Tracking" (Jan-David Wiederstein).
% Eigenständig einbindbar: % Kapitel "Data Pipeline, Baseline und Experiment-Tracking" (Jan-David Wiederstein).
% Eigenständig einbindbar: % Kapitel "Data Pipeline, Baseline und Experiment-Tracking" (Jan-David Wiederstein).
% Eigenständig einbindbar: \input{kapitel_pipeline} – benötigt tikz (arrows.meta, positioning,
% fit, backgrounds) und booktabs in der Präambel.

\section{Data Pipeline, Baseline und Experiment-Tracking}
\label{sec:pipeline}
{\small\textit{Jan-David Wiederstein}}

Aus den Befunden des Data Understanding ist eine gemeinsame Pipeline entstanden. Sie liefert
allen Ansätzen dieselben bereinigten Daten, denselben Split, dieselbe Vorverarbeitung und
dieselben Metriken -- Ergebnisse unterschiedlicher Modelle sind damit direkt vergleichbar.
Das Modell selbst baut jeder Ansatz eigenständig (Abb.~\ref{fig:pipeline}).

\subsection*{Was die Pipeline macht}

\begin{itemize}
  \item \textbf{Datenaufbereitung} (\texttt{make data}): Rohdaten werden per Schema validiert,
        zwei exakt doppelte Zeilen entfernt (sonst könnte dieselbe Messung in Training
        \emph{und} Validierung landen). Daraus entstehen ein fester 70/30-Holdout und fünf
        CV-Folds im Trainingsteil, stratifiziert nach \texttt{Crop|Stage}, damit auch seltene
        Paare wie \texttt{cotton|Harvest} (11 Zeilen) in beiden Teilen vorkommen -- per festem
        Seed bei allen identisch.
  \item \textbf{Vorverarbeitung} als sklearn-\texttt{Pipeline}, die nur auf dem Trainingsteil
        gefittet wird: 67 leere Bänder entfernen (131 bleiben), Lücken entlang der Wellenlänge
        interpolieren, \texttt{AEZ}/\texttt{Month} sowie Skalierung und Vegetationsindizes
        optional zuschaltbar.
  \item \textbf{Bewertung}: Balanced Accuracy, Macro-F1 und Samples-F1 für Crop, Stage und
        kombiniert sowie der Anteil ungültiger Crop/Stage-Paare.
\end{itemize}

\subsection*{Verwendung}

Hyperparameter werden mit \texttt{tune()} ausschließlich per Cross-Validation auf dem
Trainingsteil gesucht; das beste Modell wird mit \texttt{run()} genau einmal auf dem Holdout
bewertet. So wird nicht unbemerkt auf den Validierungsdaten optimiert.

\smallskip
\noindent\texttt{\small tuned = tune(model, TuneConfig(name=..., approach=..., param\_grid=\{...\}))}\\
\texttt{\small result = run(tuned.best\_model, RunConfig(name=..., tuning\_run=tuned.run\_id))}

\subsection*{Baseline}

\begin{center}
\small
\begin{tabular}{lcccc}
  \toprule
  Modell & BAcc Crop & BAcc Stage & BAcc kombiniert & Samples-F1 \\
  \midrule
  Dummy (häufigste Klasse) & 0{,}200 & 0{,}167 & 0{,}044 & 0{,}326 \\
  Random Forest (300 Bäume, balanced) & 0{,}890 & 0{,}870 & \textbf{0{,}743} & 0{,}878 \\
  \bottomrule
\end{tabular}
\end{center}

Der Random Forest liegt mit 0{,}743 im Streubereich seines CV-Scores ($0{,}750 \pm 0{,}019$),
überschätzt sich also nicht. Verwechselt werden vor allem Mais/Soja und benachbarte Stadien;
0{,}3\,\% der Vorhersagen sind unmögliche Paare. Jeder weitere Ansatz muss 0{,}743 schlagen.

\subsection*{Wie MLflow uns hilft}

Jeder Aufruf von \texttt{tune()} und \texttt{run()} wird automatisch protokolliert: Parameter,
Metriken, Confusion Matrices, trainiertes Modell, Hash der Daten und Git-Stand. So lassen sich
Läufe in der Oberfläche (\texttt{make mlflow}) nach Ansatz gruppieren und vergleichen, jedes
Ergebnis ist auf Code und Daten zurückführbar, und gespeicherte Modelle können direkt für die
Testvorhersage geladen werden. Wichtige Läufe kommen mit Run-ID ins Experiment-Log.

\subsection*{Was noch fehlt}

\begin{itemize}
  \item Weitere Vorverarbeitung als Optionen: Glättung (z.\,B.\ Savitzky-Golay), PCA oder
        Bandauswahl, zyklische Kodierung von \texttt{Month}, Umgang mit Ausreißern.
  \item Nur gültige Crop/Stage-Kombinationen vorhersagen (z.\,B.\ hierarchisch); gruppierter
        Split, falls sich räumliche Cluster bestätigen; Vorhersage für \texttt{test.csv}.
  \item MLflow läuft lokal (jede:r sieht nur eigene Läufe) -- ein gemeinsamer Server würde
        nur das Tracking-Modul betreffen.
\end{itemize}

\begin{figure}[p]
\centering
\begin{tikzpicture}[
    node distance=7mm and 6mm,
    >={Stealth[length=2mm]},
    box/.style={draw, rounded corners=2pt, align=center, font=\small,
                minimum width=44mm, minimum height=10mm, fill=white},
    file/.style={box, fill=gray!8},
    lbl/.style={font=\scriptsize\itshape, midway, right=1mm},
    group/.style={draw=gray!60, dashed, rounded corners=4pt, inner sep=3mm},
    gtitle/.style={font=\footnotesize\bfseries, text=gray!70!black, anchor=south west},
  ]
  \node[file] (raw) {\texttt{data/raw/train.csv}\\ Rohdaten (read-only)};
  \node[box, below=of raw] (clean) {Schema-Validierung\\ 2 Duplikate entfernen};
  \node[box, below=of clean] (split) {Stratifizierter Split nach \texttt{Crop|Stage}\\
                                       70\,\% Train / 30\,\% Holdout, 5 CV-Folds};
  \node[file, below=of split] (art) {\texttt{train\_clean.parquet}\\ \texttt{split.csv}};

  \node[box, right=42mm of raw] (pre) {Vorverarbeitung\\ leere Bänder raus, Interpolation,\\
                                        opt.\ AEZ/Month, Skalierung, Indizes};
  \node[box, below=of pre] (model) {Modell des Ansatzes\\ (z.\,B.\ Random Forest)};
  \node[box, below=of model] (tune) {\texttt{tune()}: CV auf Train};
  \node[box, below=of tune] (run) {\texttt{run()}: einmal auf Holdout\\ \texttt{evaluate()} $\to$ Metriken};

  \node[file, below=17mm of run] (mlflow) {MLflow: Parameter, Metriken,\\
                                           Modell, Daten-Hash, Git-Stand};
  \node[file, left=22mm of mlflow] (log) {Experiment-Log\\ (Doku, mit Run-ID)};

  \draw[->] (raw) -- (clean);
  \draw[->] (clean) -- (split);
  \draw[->] (split) -- (art);
  \draw[->] (art.east) -- ++(18mm,0) |- node[pos=0.25, right, font=\scriptsize\itshape]
            {\texttt{load\_split()}} (pre.west);
  \draw[->] (pre) -- (model);
  \draw[->] (model) -- (tune);
  \draw[->] (tune) -- node[lbl] {bestes Modell} (run);
  \draw[->] (run) -- node[lbl] {\texttt{log\_run()}} (mlflow);
  \draw[->] (tune.east) -- ++(6mm,0) |- (mlflow.east);
  \draw[->] (mlflow) -- node[above, font=\scriptsize\itshape] {Run-ID} (log);

  \begin{scope}[on background layer]
    \node[group, fit=(raw)(split)(art)] (g1) {};
    \node[group, fit=(pre)(run)] (g2) {};
    \node[group, fit=(mlflow)] (g3) {};
  \end{scope}
  \node[gtitle] at (g1.north west) {Einmalig: \texttt{make data}};
  \node[gtitle] at (g2.north west) {Je Ansatz: sklearn-Pipeline};
  \node[gtitle] at (g3.north west) {Tracking};
\end{tikzpicture}
\caption{Gemeinsame Data Pipeline: Die Datenaufbereitung läuft einmal und ist für alle gleich;
  Vorverarbeitung und Modell werden je Ansatz nur auf dem Trainingsteil gefittet und jeder Lauf
  in MLflow protokolliert.}
\label{fig:pipeline}
\end{figure}
 – benötigt tikz (arrows.meta, positioning,
% fit, backgrounds) und booktabs in der Präambel.

\section{Data Pipeline, Baseline und Experiment-Tracking}
\label{sec:pipeline}
{\small\textit{Jan-David Wiederstein}}

Aus den Befunden des Data Understanding ist eine gemeinsame Pipeline entstanden. Sie liefert
allen Ansätzen dieselben bereinigten Daten, denselben Split, dieselbe Vorverarbeitung und
dieselben Metriken -- Ergebnisse unterschiedlicher Modelle sind damit direkt vergleichbar.
Das Modell selbst baut jeder Ansatz eigenständig (Abb.~\ref{fig:pipeline}).

\subsection*{Was die Pipeline macht}

\begin{itemize}
  \item \textbf{Datenaufbereitung} (\texttt{make data}): Rohdaten werden per Schema validiert,
        zwei exakt doppelte Zeilen entfernt (sonst könnte dieselbe Messung in Training
        \emph{und} Validierung landen). Daraus entstehen ein fester 70/30-Holdout und fünf
        CV-Folds im Trainingsteil, stratifiziert nach \texttt{Crop|Stage}, damit auch seltene
        Paare wie \texttt{cotton|Harvest} (11 Zeilen) in beiden Teilen vorkommen -- per festem
        Seed bei allen identisch.
  \item \textbf{Vorverarbeitung} als sklearn-\texttt{Pipeline}, die nur auf dem Trainingsteil
        gefittet wird: 67 leere Bänder entfernen (131 bleiben), Lücken entlang der Wellenlänge
        interpolieren, \texttt{AEZ}/\texttt{Month} sowie Skalierung und Vegetationsindizes
        optional zuschaltbar.
  \item \textbf{Bewertung}: Balanced Accuracy, Macro-F1 und Samples-F1 für Crop, Stage und
        kombiniert sowie der Anteil ungültiger Crop/Stage-Paare.
\end{itemize}

\subsection*{Verwendung}

Hyperparameter werden mit \texttt{tune()} ausschließlich per Cross-Validation auf dem
Trainingsteil gesucht; das beste Modell wird mit \texttt{run()} genau einmal auf dem Holdout
bewertet. So wird nicht unbemerkt auf den Validierungsdaten optimiert.

\smallskip
\noindent\texttt{\small tuned = tune(model, TuneConfig(name=..., approach=..., param\_grid=\{...\}))}\\
\texttt{\small result = run(tuned.best\_model, RunConfig(name=..., tuning\_run=tuned.run\_id))}

\subsection*{Baseline}

\begin{center}
\small
\begin{tabular}{lcccc}
  \toprule
  Modell & BAcc Crop & BAcc Stage & BAcc kombiniert & Samples-F1 \\
  \midrule
  Dummy (häufigste Klasse) & 0{,}200 & 0{,}167 & 0{,}044 & 0{,}326 \\
  Random Forest (300 Bäume, balanced) & 0{,}890 & 0{,}870 & \textbf{0{,}743} & 0{,}878 \\
  \bottomrule
\end{tabular}
\end{center}

Der Random Forest liegt mit 0{,}743 im Streubereich seines CV-Scores ($0{,}750 \pm 0{,}019$),
überschätzt sich also nicht. Verwechselt werden vor allem Mais/Soja und benachbarte Stadien;
0{,}3\,\% der Vorhersagen sind unmögliche Paare. Jeder weitere Ansatz muss 0{,}743 schlagen.

\subsection*{Wie MLflow uns hilft}

Jeder Aufruf von \texttt{tune()} und \texttt{run()} wird automatisch protokolliert: Parameter,
Metriken, Confusion Matrices, trainiertes Modell, Hash der Daten und Git-Stand. So lassen sich
Läufe in der Oberfläche (\texttt{make mlflow}) nach Ansatz gruppieren und vergleichen, jedes
Ergebnis ist auf Code und Daten zurückführbar, und gespeicherte Modelle können direkt für die
Testvorhersage geladen werden. Wichtige Läufe kommen mit Run-ID ins Experiment-Log.

\subsection*{Was noch fehlt}

\begin{itemize}
  \item Weitere Vorverarbeitung als Optionen: Glättung (z.\,B.\ Savitzky-Golay), PCA oder
        Bandauswahl, zyklische Kodierung von \texttt{Month}, Umgang mit Ausreißern.
  \item Nur gültige Crop/Stage-Kombinationen vorhersagen (z.\,B.\ hierarchisch); gruppierter
        Split, falls sich räumliche Cluster bestätigen; Vorhersage für \texttt{test.csv}.
  \item MLflow läuft lokal (jede:r sieht nur eigene Läufe) -- ein gemeinsamer Server würde
        nur das Tracking-Modul betreffen.
\end{itemize}

\begin{figure}[p]
\centering
\begin{tikzpicture}[
    node distance=7mm and 6mm,
    >={Stealth[length=2mm]},
    box/.style={draw, rounded corners=2pt, align=center, font=\small,
                minimum width=44mm, minimum height=10mm, fill=white},
    file/.style={box, fill=gray!8},
    lbl/.style={font=\scriptsize\itshape, midway, right=1mm},
    group/.style={draw=gray!60, dashed, rounded corners=4pt, inner sep=3mm},
    gtitle/.style={font=\footnotesize\bfseries, text=gray!70!black, anchor=south west},
  ]
  \node[file] (raw) {\texttt{data/raw/train.csv}\\ Rohdaten (read-only)};
  \node[box, below=of raw] (clean) {Schema-Validierung\\ 2 Duplikate entfernen};
  \node[box, below=of clean] (split) {Stratifizierter Split nach \texttt{Crop|Stage}\\
                                       70\,\% Train / 30\,\% Holdout, 5 CV-Folds};
  \node[file, below=of split] (art) {\texttt{train\_clean.parquet}\\ \texttt{split.csv}};

  \node[box, right=42mm of raw] (pre) {Vorverarbeitung\\ leere Bänder raus, Interpolation,\\
                                        opt.\ AEZ/Month, Skalierung, Indizes};
  \node[box, below=of pre] (model) {Modell des Ansatzes\\ (z.\,B.\ Random Forest)};
  \node[box, below=of model] (tune) {\texttt{tune()}: CV auf Train};
  \node[box, below=of tune] (run) {\texttt{run()}: einmal auf Holdout\\ \texttt{evaluate()} $\to$ Metriken};

  \node[file, below=17mm of run] (mlflow) {MLflow: Parameter, Metriken,\\
                                           Modell, Daten-Hash, Git-Stand};
  \node[file, left=22mm of mlflow] (log) {Experiment-Log\\ (Doku, mit Run-ID)};

  \draw[->] (raw) -- (clean);
  \draw[->] (clean) -- (split);
  \draw[->] (split) -- (art);
  \draw[->] (art.east) -- ++(18mm,0) |- node[pos=0.25, right, font=\scriptsize\itshape]
            {\texttt{load\_split()}} (pre.west);
  \draw[->] (pre) -- (model);
  \draw[->] (model) -- (tune);
  \draw[->] (tune) -- node[lbl] {bestes Modell} (run);
  \draw[->] (run) -- node[lbl] {\texttt{log\_run()}} (mlflow);
  \draw[->] (tune.east) -- ++(6mm,0) |- (mlflow.east);
  \draw[->] (mlflow) -- node[above, font=\scriptsize\itshape] {Run-ID} (log);

  \begin{scope}[on background layer]
    \node[group, fit=(raw)(split)(art)] (g1) {};
    \node[group, fit=(pre)(run)] (g2) {};
    \node[group, fit=(mlflow)] (g3) {};
  \end{scope}
  \node[gtitle] at (g1.north west) {Einmalig: \texttt{make data}};
  \node[gtitle] at (g2.north west) {Je Ansatz: sklearn-Pipeline};
  \node[gtitle] at (g3.north west) {Tracking};
\end{tikzpicture}
\caption{Gemeinsame Data Pipeline: Die Datenaufbereitung läuft einmal und ist für alle gleich;
  Vorverarbeitung und Modell werden je Ansatz nur auf dem Trainingsteil gefittet und jeder Lauf
  in MLflow protokolliert.}
\label{fig:pipeline}
\end{figure}
 – benötigt tikz (arrows.meta, positioning,
% fit, backgrounds) und booktabs in der Präambel.

\section{Data Pipeline, Baseline und Experiment-Tracking}
\label{sec:pipeline}
{\small\textit{Jan-David Wiederstein}}

Aus den Befunden des Data Understanding ist eine gemeinsame Pipeline entstanden. Sie liefert
allen Ansätzen dieselben bereinigten Daten, denselben Split, dieselbe Vorverarbeitung und
dieselben Metriken -- Ergebnisse unterschiedlicher Modelle sind damit direkt vergleichbar.
Das Modell selbst baut jeder Ansatz eigenständig (Abb.~\ref{fig:pipeline}).

\subsection*{Was die Pipeline macht}

\begin{itemize}
  \item \textbf{Datenaufbereitung} (\texttt{make data}): Rohdaten werden per Schema validiert,
        zwei exakt doppelte Zeilen entfernt (sonst könnte dieselbe Messung in Training
        \emph{und} Validierung landen). Daraus entstehen ein fester 70/30-Holdout und fünf
        CV-Folds im Trainingsteil, stratifiziert nach \texttt{Crop|Stage}, damit auch seltene
        Paare wie \texttt{cotton|Harvest} (11 Zeilen) in beiden Teilen vorkommen -- per festem
        Seed bei allen identisch.
  \item \textbf{Vorverarbeitung} als sklearn-\texttt{Pipeline}, die nur auf dem Trainingsteil
        gefittet wird: 67 leere Bänder entfernen (131 bleiben), Lücken entlang der Wellenlänge
        interpolieren, \texttt{AEZ}/\texttt{Month} sowie Skalierung und Vegetationsindizes
        optional zuschaltbar.
  \item \textbf{Bewertung}: Balanced Accuracy, Macro-F1 und Samples-F1 für Crop, Stage und
        kombiniert sowie der Anteil ungültiger Crop/Stage-Paare.
\end{itemize}

\subsection*{Verwendung}

Hyperparameter werden mit \texttt{tune()} ausschließlich per Cross-Validation auf dem
Trainingsteil gesucht; das beste Modell wird mit \texttt{run()} genau einmal auf dem Holdout
bewertet. So wird nicht unbemerkt auf den Validierungsdaten optimiert.

\smallskip
\noindent\texttt{\small tuned = tune(model, TuneConfig(name=..., approach=..., param\_grid=\{...\}))}\\
\texttt{\small result = run(tuned.best\_model, RunConfig(name=..., tuning\_run=tuned.run\_id))}

\subsection*{Baseline}

\begin{center}
\small
\begin{tabular}{lcccc}
  \toprule
  Modell & BAcc Crop & BAcc Stage & BAcc kombiniert & Samples-F1 \\
  \midrule
  Dummy (häufigste Klasse) & 0{,}200 & 0{,}167 & 0{,}044 & 0{,}326 \\
  Random Forest (300 Bäume, balanced) & 0{,}890 & 0{,}870 & \textbf{0{,}743} & 0{,}878 \\
  \bottomrule
\end{tabular}
\end{center}

Der Random Forest liegt mit 0{,}743 im Streubereich seines CV-Scores ($0{,}750 \pm 0{,}019$),
überschätzt sich also nicht. Verwechselt werden vor allem Mais/Soja und benachbarte Stadien;
0{,}3\,\% der Vorhersagen sind unmögliche Paare. Jeder weitere Ansatz muss 0{,}743 schlagen.

\subsection*{Wie MLflow uns hilft}

Jeder Aufruf von \texttt{tune()} und \texttt{run()} wird automatisch protokolliert: Parameter,
Metriken, Confusion Matrices, trainiertes Modell, Hash der Daten und Git-Stand. So lassen sich
Läufe in der Oberfläche (\texttt{make mlflow}) nach Ansatz gruppieren und vergleichen, jedes
Ergebnis ist auf Code und Daten zurückführbar, und gespeicherte Modelle können direkt für die
Testvorhersage geladen werden. Wichtige Läufe kommen mit Run-ID ins Experiment-Log.

\subsection*{Was noch fehlt}

\begin{itemize}
  \item Weitere Vorverarbeitung als Optionen: Glättung (z.\,B.\ Savitzky-Golay), PCA oder
        Bandauswahl, zyklische Kodierung von \texttt{Month}, Umgang mit Ausreißern.
  \item Nur gültige Crop/Stage-Kombinationen vorhersagen (z.\,B.\ hierarchisch); gruppierter
        Split, falls sich räumliche Cluster bestätigen; Vorhersage für \texttt{test.csv}.
  \item MLflow läuft lokal (jede:r sieht nur eigene Läufe) -- ein gemeinsamer Server würde
        nur das Tracking-Modul betreffen.
\end{itemize}

\begin{figure}[p]
\centering
\begin{tikzpicture}[
    node distance=7mm and 6mm,
    >={Stealth[length=2mm]},
    box/.style={draw, rounded corners=2pt, align=center, font=\small,
                minimum width=44mm, minimum height=10mm, fill=white},
    file/.style={box, fill=gray!8},
    lbl/.style={font=\scriptsize\itshape, midway, right=1mm},
    group/.style={draw=gray!60, dashed, rounded corners=4pt, inner sep=3mm},
    gtitle/.style={font=\footnotesize\bfseries, text=gray!70!black, anchor=south west},
  ]
  \node[file] (raw) {\texttt{data/raw/train.csv}\\ Rohdaten (read-only)};
  \node[box, below=of raw] (clean) {Schema-Validierung\\ 2 Duplikate entfernen};
  \node[box, below=of clean] (split) {Stratifizierter Split nach \texttt{Crop|Stage}\\
                                       70\,\% Train / 30\,\% Holdout, 5 CV-Folds};
  \node[file, below=of split] (art) {\texttt{train\_clean.parquet}\\ \texttt{split.csv}};

  \node[box, right=42mm of raw] (pre) {Vorverarbeitung\\ leere Bänder raus, Interpolation,\\
                                        opt.\ AEZ/Month, Skalierung, Indizes};
  \node[box, below=of pre] (model) {Modell des Ansatzes\\ (z.\,B.\ Random Forest)};
  \node[box, below=of model] (tune) {\texttt{tune()}: CV auf Train};
  \node[box, below=of tune] (run) {\texttt{run()}: einmal auf Holdout\\ \texttt{evaluate()} $\to$ Metriken};

  \node[file, below=17mm of run] (mlflow) {MLflow: Parameter, Metriken,\\
                                           Modell, Daten-Hash, Git-Stand};
  \node[file, left=22mm of mlflow] (log) {Experiment-Log\\ (Doku, mit Run-ID)};

  \draw[->] (raw) -- (clean);
  \draw[->] (clean) -- (split);
  \draw[->] (split) -- (art);
  \draw[->] (art.east) -- ++(18mm,0) |- node[pos=0.25, right, font=\scriptsize\itshape]
            {\texttt{load\_split()}} (pre.west);
  \draw[->] (pre) -- (model);
  \draw[->] (model) -- (tune);
  \draw[->] (tune) -- node[lbl] {bestes Modell} (run);
  \draw[->] (run) -- node[lbl] {\texttt{log\_run()}} (mlflow);
  \draw[->] (tune.east) -- ++(6mm,0) |- (mlflow.east);
  \draw[->] (mlflow) -- node[above, font=\scriptsize\itshape] {Run-ID} (log);

  \begin{scope}[on background layer]
    \node[group, fit=(raw)(split)(art)] (g1) {};
    \node[group, fit=(pre)(run)] (g2) {};
    \node[group, fit=(mlflow)] (g3) {};
  \end{scope}
  \node[gtitle] at (g1.north west) {Einmalig: \texttt{make data}};
  \node[gtitle] at (g2.north west) {Je Ansatz: sklearn-Pipeline};
  \node[gtitle] at (g3.north west) {Tracking};
\end{tikzpicture}
\caption{Gemeinsame Data Pipeline: Die Datenaufbereitung läuft einmal und ist für alle gleich;
  Vorverarbeitung und Modell werden je Ansatz nur auf dem Trainingsteil gefittet und jeder Lauf
  in MLflow protokolliert.}
\label{fig:pipeline}
\end{figure}
 – benötigt tikz (arrows.meta, positioning,
% fit, backgrounds) und booktabs in der Präambel.

\section{Data Pipeline, Baseline und Experiment-Tracking}
\label{sec:pipeline}
{\small\textit{Jan-David Wiederstein}}

Aus den Befunden des Data Understanding ist eine gemeinsame Pipeline entstanden. Sie liefert
allen Ansätzen dieselben bereinigten Daten, denselben Split, dieselbe Vorverarbeitung und
dieselben Metriken -- Ergebnisse unterschiedlicher Modelle sind damit direkt vergleichbar.
Das Modell selbst baut jeder Ansatz eigenständig (Abb.~\ref{fig:pipeline}).

\subsection*{Was die Pipeline macht}

\begin{itemize}
  \item \textbf{Datenaufbereitung} (\texttt{make data}): Rohdaten werden per Schema validiert,
        zwei exakt doppelte Zeilen entfernt (sonst könnte dieselbe Messung in Training
        \emph{und} Validierung landen). Daraus entstehen ein fester 70/30-Holdout und fünf
        CV-Folds im Trainingsteil, stratifiziert nach \texttt{Crop|Stage}, damit auch seltene
        Paare wie \texttt{cotton|Harvest} (11 Zeilen) in beiden Teilen vorkommen -- per festem
        Seed bei allen identisch.
  \item \textbf{Vorverarbeitung} als sklearn-\texttt{Pipeline}, die nur auf dem Trainingsteil
        gefittet wird: 67 leere Bänder entfernen (131 bleiben), Lücken entlang der Wellenlänge
        interpolieren, \texttt{AEZ}/\texttt{Month} sowie Skalierung und Vegetationsindizes
        optional zuschaltbar.
  \item \textbf{Bewertung}: Balanced Accuracy, Macro-F1 und Samples-F1 für Crop, Stage und
        kombiniert sowie der Anteil ungültiger Crop/Stage-Paare.
\end{itemize}

\subsection*{Verwendung}

Hyperparameter werden mit \texttt{tune()} ausschließlich per Cross-Validation auf dem
Trainingsteil gesucht; das beste Modell wird mit \texttt{run()} genau einmal auf dem Holdout
bewertet. So wird nicht unbemerkt auf den Validierungsdaten optimiert.

\smallskip
\noindent\texttt{\small tuned = tune(model, TuneConfig(name=..., approach=..., param\_grid=\{...\}))}\\
\texttt{\small result = run(tuned.best\_model, RunConfig(name=..., tuning\_run=tuned.run\_id))}

\subsection*{Baseline}

\begin{center}
\small
\begin{tabular}{lcccc}
  \toprule
  Modell & BAcc Crop & BAcc Stage & BAcc kombiniert & Samples-F1 \\
  \midrule
  Dummy (häufigste Klasse) & 0{,}200 & 0{,}167 & 0{,}044 & 0{,}326 \\
  Random Forest (300 Bäume, balanced) & 0{,}890 & 0{,}870 & \textbf{0{,}743} & 0{,}878 \\
  \bottomrule
\end{tabular}
\end{center}

Der Random Forest liegt mit 0{,}743 im Streubereich seines CV-Scores ($0{,}750 \pm 0{,}019$),
überschätzt sich also nicht. Verwechselt werden vor allem Mais/Soja und benachbarte Stadien;
0{,}3\,\% der Vorhersagen sind unmögliche Paare. Jeder weitere Ansatz muss 0{,}743 schlagen.

\subsection*{Wie MLflow uns hilft}

Jeder Aufruf von \texttt{tune()} und \texttt{run()} wird automatisch protokolliert: Parameter,
Metriken, Confusion Matrices, trainiertes Modell, Hash der Daten und Git-Stand. So lassen sich
Läufe in der Oberfläche (\texttt{make mlflow}) nach Ansatz gruppieren und vergleichen, jedes
Ergebnis ist auf Code und Daten zurückführbar, und gespeicherte Modelle können direkt für die
Testvorhersage geladen werden. Wichtige Läufe kommen mit Run-ID ins Experiment-Log.

\subsection*{Was noch fehlt}

\begin{itemize}
  \item Weitere Vorverarbeitung als Optionen: Glättung (z.\,B.\ Savitzky-Golay), PCA oder
        Bandauswahl, zyklische Kodierung von \texttt{Month}, Umgang mit Ausreißern.
  \item Nur gültige Crop/Stage-Kombinationen vorhersagen (z.\,B.\ hierarchisch); gruppierter
        Split, falls sich räumliche Cluster bestätigen; Vorhersage für \texttt{test.csv}.
  \item MLflow läuft lokal (jede:r sieht nur eigene Läufe) -- ein gemeinsamer Server würde
        nur das Tracking-Modul betreffen.
\end{itemize}

\begin{figure}[p]
\centering
\begin{tikzpicture}[
    node distance=7mm and 6mm,
    >={Stealth[length=2mm]},
    box/.style={draw, rounded corners=2pt, align=center, font=\small,
                minimum width=44mm, minimum height=10mm, fill=white},
    file/.style={box, fill=gray!8},
    lbl/.style={font=\scriptsize\itshape, midway, right=1mm},
    group/.style={draw=gray!60, dashed, rounded corners=4pt, inner sep=3mm},
    gtitle/.style={font=\footnotesize\bfseries, text=gray!70!black, anchor=south west},
  ]
  \node[file] (raw) {\texttt{data/raw/train.csv}\\ Rohdaten (read-only)};
  \node[box, below=of raw] (clean) {Schema-Validierung\\ 2 Duplikate entfernen};
  \node[box, below=of clean] (split) {Stratifizierter Split nach \texttt{Crop|Stage}\\
                                       70\,\% Train / 30\,\% Holdout, 5 CV-Folds};
  \node[file, below=of split] (art) {\texttt{train\_clean.parquet}\\ \texttt{split.csv}};

  \node[box, right=42mm of raw] (pre) {Vorverarbeitung\\ leere Bänder raus, Interpolation,\\
                                        opt.\ AEZ/Month, Skalierung, Indizes};
  \node[box, below=of pre] (model) {Modell des Ansatzes\\ (z.\,B.\ Random Forest)};
  \node[box, below=of model] (tune) {\texttt{tune()}: CV auf Train};
  \node[box, below=of tune] (run) {\texttt{run()}: einmal auf Holdout\\ \texttt{evaluate()} $\to$ Metriken};

  \node[file, below=17mm of run] (mlflow) {MLflow: Parameter, Metriken,\\
                                           Modell, Daten-Hash, Git-Stand};
  \node[file, left=22mm of mlflow] (log) {Experiment-Log\\ (Doku, mit Run-ID)};

  \draw[->] (raw) -- (clean);
  \draw[->] (clean) -- (split);
  \draw[->] (split) -- (art);
  \draw[->] (art.east) -- ++(18mm,0) |- node[pos=0.25, right, font=\scriptsize\itshape]
            {\texttt{load\_split()}} (pre.west);
  \draw[->] (pre) -- (model);
  \draw[->] (model) -- (tune);
  \draw[->] (tune) -- node[lbl] {bestes Modell} (run);
  \draw[->] (run) -- node[lbl] {\texttt{log\_run()}} (mlflow);
  \draw[->] (tune.east) -- ++(6mm,0) |- (mlflow.east);
  \draw[->] (mlflow) -- node[above, font=\scriptsize\itshape] {Run-ID} (log);

  \begin{scope}[on background layer]
    \node[group, fit=(raw)(split)(art)] (g1) {};
    \node[group, fit=(pre)(run)] (g2) {};
    \node[group, fit=(mlflow)] (g3) {};
  \end{scope}
  \node[gtitle] at (g1.north west) {Einmalig: \texttt{make data}};
  \node[gtitle] at (g2.north west) {Je Ansatz: sklearn-Pipeline};
  \node[gtitle] at (g3.north west) {Tracking};
\end{tikzpicture}
\caption{Gemeinsame Data Pipeline: Die Datenaufbereitung läuft einmal und ist für alle gleich;
  Vorverarbeitung und Modell werden je Ansatz nur auf dem Trainingsteil gefittet und jeder Lauf
  in MLflow protokolliert.}
\label{fig:pipeline}
\end{figure}
